% META: {"id":"1SPE-PROBCOND-CO-048","chapitre":"1SPE-PROBA-COND","type_objet":"corrige","exercice_ref":"1SPE-PROBCOND-EX-048","capacites_codes":["C5"],"sources_inspiration":[],"status":"approved","origin":"generated","status_history":["generated","audited","accepted","approved"],"release_acceptance":"RELEASE_OWNER_BATCH_ACCEPTANCE","acceptance_closure_digest":"sha256:9b3ccf9a81c5520fbb7e03b7b2d3e7bf2057a02834b6d908bf3be5f49f3b553f","acceptance_receipt_digest":"sha256:4fad86f0282e0fa4e0419cca1ad6379ae7af5a280c04a4a90880f90a1c044242"}
% BEGIN-VERIFY
% from sympy import Rational
% assert Rational(9, 100) == Rational(4,5)*Rational(1,20)+Rational(1,5)*Rational(1,4)
% assert Rational(5, 9) == (Rational(1,5)*Rational(1,4))/Rational(9,100)
% END-VERIFY
\begin{corrige}{1SPE-PROBCOND-EX-048}

\textbf{1.} $P(A) = \frac{4}{5}\times\frac{1}{20}+\frac{1}{5}\times\frac{1}{4} = \frac{4}{100}+\frac{5}{100} = \frac{9}{100}$.

\textbf{2.} $P_A(R) = \frac{P(R \cap A)}{P(A)} = \frac{1/20}{9/100} = \frac{100}{180} = \frac{5}{9} \approx 55{,}6\,\%$.

\textbf{3.} Plus de la moitie des accidentes sont des conducteurs a risque, alors qu'ils ne representent que $20\,\%$ de la population.

\end{corrige}
